\section{Future Work}
Due to computational constraints, our post-training experiments focus primarily on mathematical and code reasoning tasks. We do not extensively explore other important capabilities studied in prior work, such as scientific question answering, instruction following, alignment with human conversational preferences, or agentic task-solving. A natural direction for future work is to continue to post-train Loopie on a broader set of tasks and preference signals, which may further improve its practical utility and general-purpose capabilities. Also, due to limited computational resources, we were unable to conduct a sufficiently comprehensive ablation study of supervised pre-training, which we leave for future work.

In addition, our study focuses primarily on matching compute budgets during pre-training, and we have not yet conducted systematic studies of inference-time computation. Matching and optimizing inference-time compute remain important directions for future work; prior work such as the Parallel Loop Transformer~\citep{wu2025parallellooptransformerefficient} provides a promising example of this approach.

Finally, our study intentionally focuses on applying looped Transformers to a relatively clean and well-controlled base architecture, Qwen3-30B-A3B, in order to avoid confounding factors introduced by architectural variations. As a result, we do not investigate several recent architectural advances that may be complementary to our method. Exploring how Loopie interacts with these newer designs remains an important direction for future work.

\section{Conclusion}
We introduced Loopie, a family of looped MoE language models that makes recurrent depth competitive under a matched pre-training compute budget. By combining layer-loop recurrence with a hardware-aware scaling recipe, Loopie consistently outperforms compute-matched vanilla Transformer baselines across model scales. A large-scale post-training pipeline based on Supervised Pre-training and reinforcement learning further equips Loopie with strong mathematical reasoning and coding abilities. These results suggest that recurrent computation, when jointly optimized with architecture and training efficiency, can serve as a practical scaling axis for large language models.
